\textbf{\textit{Abstract}} - Este artículo presenta el desarrollo y la evaluación de un entorno educativo gamificado en Roblox, diseñado específicamente para la enseñanza práctica de metodologías de aseguramiento de calidad de software (QA). El sistema, denominado "Shop \& Jump Simulator (Bug Edition)", implementa una arquitectura modular con errores lógicos intencionales en sus subsistemas de economía, progresión y sincronización multijugador. El objetivo principal es proporcionar a los estudiantes un escenario controlado donde puedan aplicar niveles de prueba manual, unitaria, de integración y automatizada mediante el uso de TestService. Los resultados demuestran que la exposición a fallos distribuidos y flujos de integración rotos fortalece la capacidad diagnóstica del alumno, permitiendo identificar vulnerabilidades críticas como saldos negativos, inconsistencias en multiplicadores de recompensa y falta de validación en la comunicación cliente-servidor.

Todos los códigos fuente, se encuentran dentro del repositorio del proyecto, el cual se encuentra disponible en GitHub: \url{https://github.com/ArelyX1/TestingUCSM/tree/master/F1/Teo1}.  
\vspace{0.5em}

\textbf{\textit{Index Terms}} - Software Testing, Unit Testing, Integration Testing, Test Automation, Gamification, VS Code